\chapter{Introducción}
\label{ch:introduccion}

Este informe documenta el desarrollo y la resolución del Ejercicio Práctico N° 3 del Trabajo Práctico N° 2.
Las actividades propuestas en esta etapa tienen como objetivo principal abandonar el uso de punteros
manuales de hardware (Bare Metal puro) y migrar la base de código existente hacia el estándar CMSIS
(Cortex Microcontroller Software Interface Standard).

El propósito general de este ejercicio es adaptar el código del Ejercicio N° 2, el cual configura
interrupciones externas, canales del ADC, y la interfaz GPIO, para utilizar las estructuras y macros
definidas en \lstinline{stm32f103xb.h}. Esto permite un desarrollo más robusto, una mejora significativa en la
legibilidad del código y una mayor facilidad de mantenimiento. Adicionalmente, se reemplaza el retardo
basado en bucles (\textit{software delay}) por una solución determinista por hardware, configurando el periférico
interno del núcleo Cortex-M3 denominado \textit{SysTick} para generar retardos precisos de \qty{1}{\ms}.
